\section{Discussion}
\label{sec:discussion}

\paragraph{A random split rewards memorisation.} Under the random protocol at most \V{full.random.maxnew} of the \V{full.split.random.s42.ntest} test tasks of a split come from users without training tasks, and \V{rf.full.dup.random.min} to \V{rf.full.dup.random.max} have an exact feature duplicate in training. The models with interactions lead there: LightGBM is \V{full.random.lgbm.absdiff} AUC above \kaam{} on the full trace, and LightGBM and the EBM with pairwise terms are \V{orig.random.lgbm.absdiff} and \V{orig.random.ebmi.absdiff} above it on the original sample. Under the user-grouped protocol that lead is gone: on the four folds without User~1, LightGBM scores \V{full.gnou.lgbm.auc} and the MLP and \kaam{} \V{full.gnou.mlp.auc} and \V{full.gnou.kaam.auc}, and no model is clearly best. This pattern is consistent with the flexible models' random-split gain being knowledge of particular users and of duplicated requests. It does not show that interactions themselves fail to transfer: LightGBM's grouped fits were cut after a few boosting rounds by early stopping on validation users, and a regularised or interaction-constrained boosting model would be needed to test that.

\paragraph{What transfers and what does not.} Two settings matter for a failure predictor in production. For users already in the training data, the features carry over to a later period only weakly: on the last \V{prot.time.test} of the trace every model scores between \V{full.timeknown.minauc} and \V{full.timeknown.maxauc} on known users, although on the validation period just before it, a burst in which User~1 submits \V{rf.full.time.val.u1} of the tasks, the same models scored between \V{full.time.minvalauc} and \V{full.time.maxvalauc}. The trace changes near its end; the failure rate is \V{full.split.time.t.ratetrain} in the training period, \V{full.split.time.t.rateval} in the validation period and \V{full.split.time.t.ratetest} in the test period. For new users the features carry a weak signal at best: the user-grouped AUCs lie between \V{full.grouped.minauc} and \V{full.grouped.maxauc} on the full trace with wide intervals, and the new users of the time-ordered test period, few in number, are scored between \V{full.timenew.minauc} and \V{full.timenew.maxauc}. A deployment could at most support an operator who reviews the recurring pipelines of known users; it should not be used to judge the first jobs of a new user, and its accuracy on a random split describes neither case. Whether frequent retraining would close the gap to a later period is not tested here.

\paragraph{What the audit contributed.} The finding that drives this paper came from reading the model. The \feat{plan_cpu} curve assigns a large rise to one-core requests; the case audit of the most confident catch puts most of its score on that single term; counting the training tasks behind it shows a request submitted almost entirely by User~1; and once User~1 is removed from the training data, one-core tasks fail about as often as other tasks. The same chain explains the most confident false alarm, a one-core task of another user that inherits User~1's term. Each step uses an exact property of the additive model: the curve is the model's whole treatment of the feature, and the contributions are exact SHAP values. Any additive model, the additive EBM included, supports the same chain; the KAN form adds smooth curves with an explicit smoothness penalty, not a different kind of explanation. A global importance score would show that \feat{plan_cpu} matters; it would not show by itself that its effect comes from one user's recurring job.

\paragraph{The smoothness penalty is chosen by the protocol.} On random splits validation preferred no penalty, and the request curves oscillate between observed values; when the validation tasks came from other users or a later period, validation preferred the strongest penalty. A smooth curve is easier to read, and validation prefers it when the validation population differs from the training population.

\paragraph{What changed from the first version.} The first version reported that \kaam{} matched a parameter-matched MLP on one random split of the original sample (\V{v1.kaam.auc} against \V{v1.mlp.auc}) and that every model fell below chance on one user-grouped partition (\kaam{} \V{v1.grouped.kaam}). With five splits, a stronger baseline set and the leak removed, \kaam{} still beats the parameter-matched MLP on the original sample (\V{orig.random.mlp.diff} AUC) but not on the full trace (\V{full.random.mlp.diff}), and it trails LightGBM on both datasets and the EBM with interactions on the original sample. With five user-grouped folds the picture on new users is no longer uniform, and no model is clearly best: the fold AUCs range from \V{full.fold.minauc} to \V{full.fold.maxauc}, and \V{full.fold.nbelow} of the \V{full.fold.npairs} fold-model pairs fall below chance. The first version's single partition was one draw from this spread. The conclusion that the features transfer poorly to new users stands, and it now rests on five folds, two datasets and six models.
